\documentclass{ibetest}
\begin{document}
\maketest{Progress Test 3.C --- Elec-S1}
         {3.C \,$\cdot$\, Operating point of a circuit}
         {10 minutes}{14}

% ---------------------------------------------------------------------------
\exercise{Maximum power}{1 mark}
A source has $\ETh = \qty{12}{V}$ and $\RTh = \qty{6.0}{\ohm}$. Give the maximum power it
can deliver to a resistive load.
\answer{1}{2.2cm}{%
  Matching $R_c = \RTh$: \quad
  $P_c^{\max} = \dfrac{\ETh^2}{4\RTh} = \dfrac{144}{4\times6.0}\ \mathrm{W}$ \qquad
  $\boxed{P_c^{\max} = 6.0\ \mathrm{W}}$}
\begin{scheme}
\pt{1}{1 pt}{$P_c^{\max} = 6.0\ \mathrm{W}$.}
\end{scheme}

% ---------------------------------------------------------------------------
\exercise{The operating point}{5 marks}
A linear network, described by its Thévenin model $(\ETh, \RTh)$, feeds a load resistor
$R_c$. Define the operating point. Write down the characteristic of the source and that of
the load, and deduce the analytical expressions of $u^\star$ and $i^\star$.
\answer{2,2,1}{7.0cm}{%
  \begin{minipage}[c]{0.6\linewidth}\raggedright
    The \textbf{operating point} is the actual state of the circuit: the unique pair
    $(u^\star, i^\star)$ that satisfies \emph{simultaneously} the characteristic of the
    source and that of the load. Graphically, it is the \textbf{intersection} of the two
    characteristics.\\[4pt]
    Source (generator convention): $u = \ETh - \RTh\, i$\\[2pt]
    Load (receiver convention): $u = R_c\, i$
  \end{minipage}\hfill
  \begin{minipage}[c]{0.37\linewidth}\centering
  \begin{tikzpicture}[scale=0.82, >={Latex[length=1.6mm]}, line width=0.5pt]
    \draw[->] (0,0) -- (3.5,0) node[below] {$u$};
    \draw[->] (0,0) -- (0,2.45) node[left] {$i$};
    \draw[thick] (3,0) -- (0,2);
    \draw[thick, dashed] (0,0) -- (3.3,1.65);
    \fill (1.714,0.857) circle (1.6pt);
    \draw[densely dotted] (1.714,0) -- (1.714,0.857) -- (0,0.857);
    \node[below, font=\scriptsize] at (3,0) {$\ETh$};
    \node[left, font=\scriptsize] at (0,2) {$\ETh/\RTh$};
    \node[below, font=\scriptsize] at (1.714,0) {$u^\star$};
    \node[left, font=\scriptsize] at (0,0.857) {$i^\star$};
    \node[font=\scriptsize, anchor=north] at (3.1,1.45) {load};
    \node[font=\scriptsize, anchor=west] at (0.55,1.85) {source};
  \end{tikzpicture}
  \end{minipage}\\[5pt]
  $R_c\, i^\star = \ETh - \RTh\, i^\star$ \quad$\Rightarrow$\quad
  $\boxed{i^\star = \dfrac{\ETh}{\RTh + R_c}}$ \qquad
  $\boxed{u^\star = R_c\, i^\star = \dfrac{R_c}{\RTh + R_c}\,\ETh}$}
\begin{scheme}
\pt{1}{2 pts}{Pair $(u^\star, i^\star)$ satisfying both characteristics (1~pt); intersection
  of the two curves (1~pt).}
\pt{2}{2 pts}{$u = \ETh - \RTh i$ (1~pt) and $u = R_c i$ (1~pt).}
\pt{3}{1 pt}{Both $i^\star$ and $u^\star$.}
\end{scheme}

% ---------------------------------------------------------------------------
\exercise{Power transfer}{3 marks}
\question{When the load resistance tends to infinity (open circuit):}
\opttwo{0}{$u^\star \to 0$ and $i^\star \to \ETh/\RTh$}{1}{$u^\star \to \ETh$ and
        $i^\star \to 0$}{0}{$u^\star \to \ETh$ and $i^\star \to \ETh/\RTh$}{0}{$u^\star \to 0$
        and $i^\star \to 0$}
\why{$i^\star = \ETh/(\RTh + R_c) \to 0$ and $u^\star \to \ETh$: the point reaches the
  voltage axis, and the load receives no power.}
\question{When the load is matched to the source ($R_c = \RTh$), the efficiency $\eta$ of the
  transfer is:}
\opts{0}{100\,\%}{0}{75\,\%}{1}{50\,\%}{0}{25\,\%}
\why{$\eta = R_c/(\RTh + R_c) = 1/2$: as much power is lost in $\RTh$ as reaches the load.}
\question{To transmit electric power over long distances, one chooses:}
\optlist{0}{matching, $R_c = \RTh$, to obtain the maximum power;}
        {1}{a high efficiency, $R_c \gg \RTh$, to limit the losses;}
        {0}{$R_c \ll \RTh$, to obtain the maximum current;}
        {0}{a short circuit at the end of the line.}
\why{energy is costly and losses cause heating, so efficiency wins. Matching is for tiny
  signals (antenna, sensor), where extracting the most power is what matters.}

% ---------------------------------------------------------------------------
\exercise{Operating point, power and efficiency}{5 marks}
A source with $\ETh = \qty{12}{V}$ and $\RTh = \qty{4.0}{\ohm}$ feeds a load
$R_c = \qty{8.0}{\ohm}$. Determine the operating point $(u^\star, i^\star)$, the power $P_c$
received by the load and the efficiency $\eta$ of the transfer. Check the power balance.
\data{$\ETh = \qty{12}{V}$ \hspace{14mm} $\RTh = \qty{4.0}{\ohm}$ \hspace{14mm}
      $R_c = \qty{8.0}{\ohm}$}
\answer{1,1,1,1,1}{6.6cm}{%
  $i^\star = \dfrac{\ETh}{\RTh + R_c} = \dfrac{12\ \mathrm{V}}{12\ \Omega}$ \qquad
  $\boxed{i^\star = 1.0\ \mathrm{A}}$ \qquad\qquad
  $u^\star = R_c\, i^\star$ \qquad $\boxed{u^\star = 8.0\ \mathrm{V}}$\\[4pt]
  $P_c = R_c\, {i^\star}^2 = 8.0\ \Omega\times(1.0\ \mathrm{A})^2$ \qquad
  $\boxed{P_c = 8.0\ \mathrm{W}}$\\[4pt]
  $\eta = \dfrac{P_c}{\ETh\, i^\star} = \dfrac{R_c}{\RTh + R_c} = \dfrac{8.0}{12}$ \qquad
  $\boxed{\eta = 2/3 \approx 67\,\%}$\\[4pt]
  Balance: produced by the ideal source $\ETh i^\star = 12\ \mathrm{W}$ $=$ lost in $\RTh$,
  $\RTh {i^\star}^2 = 4.0\ \mathrm{W}$, $+$ received by the load, $8.0\ \mathrm{W}$.
  \checkmark}
\begin{scheme}
\pt{1}{1 pt}{$i^\star = 1.0\ \mathrm{A}$.}
\pt{2}{1 pt}{$u^\star = 8.0\ \mathrm{V}$.}
\pt{3}{1 pt}{$P_c = 8.0\ \mathrm{W}$.}
\pt{4}{1 pt}{$\eta = 2/3 \approx 67\,\%$.}
\pt{5}{1 pt}{Balance $12\ \mathrm{W} = 4.0\ \mathrm{W} + 8.0\ \mathrm{W}$.}
\end{scheme}
\begin{tip}
$\eta$ is the ratio of the power received by the load to the power produced by the
\emph{ideal} source $\ETh i^\star$, not to $P_c^{\max}$.
\end{tip}

\finish
\end{document}
